% GENERATED FILE. Edit canonical lessons, then run npm run generate:books.
% canonical-source-hash: fnv1a64:62e89de571bdf568
% canonical-lessons: IT-C489-ammalarsi, IT-C489-sintomo, IT-C489-diagnosi, IT-C489-infezione, IT-C489-guarire

\chapter{Ammalarsi e guarire}
\label{ch:it-ammalarsi-e-guarire}
\hlchaptermodality{489}

\begin{chapteropening}
\textbf{By the end of this chapter:} \emph{I can describe getting ill: symptoms, diagnosis, infection and recovery.}

Complete the last lesson of chapter 489: I can describe getting ill: symptoms, diagnosis, infection and recovery.
\end{chapteropening}

\section[ammalarsi]{ammalarsi — to fall ill}

\label{lesson:IT-C489-ammalarsi}



\begin{quote}
Before the new one: say the Italian for to operate, then the Italian for anaesthesia.
\end{quote}

\subsection*{You'll want to know: ammalarsi}
\textbf{ammalarsi} — \textquotedblleft{}to fall ill\textquotedblright{}.

\begin{quote}
Ogni inverno mi ammalo almeno una volta.
\end{quote}

\begin{quote}
\emph{Every winter I fall ill at least once.}
\end{quote}

It takes \emph{essere}: \emph{mi sono ammalato}. Built on \emph{malato} \textquotedblleft{}ill\textquotedblright{}.

\begin{cousinweb}
\emph{ammalarsi} is built on \emph{malato} \textquotedblleft{}ill\textquotedblright{}, usually traced to Latin \emph{male habitus} \textquotedblleft{}in poor condition\textquotedblright{} (Lewis \& Short).
\end{cousinweb}

\begin{grammarlens}[title={What you've built}]
One more word for telling, explaining and describing at B1.
\end{grammarlens}

\subsection*{Your turn}
\begin{itemize}
\raggedright
  \item \emph{Say it:} \emph{ammalarsi}
  \item \emph{Say it:} \emph{ammalarsi}, once more
  \item \emph{Say it:} a sentence of your own with \emph{ammalarsi}
  \item \emph{Recall:} say the Italian for to operate, then the Italian for anaesthesia, then say \emph{ammalarsi} again
\end{itemize}

\begin{tcolorbox}[breakable,colback=teal!4,colframe=teal!35!black,title={Before you move on}]
What does ammalarsi mean? (\textquotedblleft{}To fall ill\textquotedblright{}.) Say it once more.
\end{tcolorbox}
\section[il sintomo]{il sintomo — symptom}

\label{lesson:IT-C489-sintomo}



\begin{quote}
Before the new one: say the Italian for anaesthesia, then the Italian for to fall ill.
\end{quote}

\subsection*{You'll want to know: il sintomo}
\textbf{il sintomo} — \textquotedblleft{}symptom\textquotedblright{}.

\begin{quote}
Quali sono i sintomi dell'influenza?
\end{quote}

\begin{quote}
\emph{What are the symptoms of flu?}
\end{quote}

The plural is \emph{i sintomi}.

\begin{cousinweb}
\emph{sintomo} is a learned borrowing, through Late Latin \emph{symptōma}, from Greek \emph{sýmptōma} \textquotedblleft{}a happening, a mishap\textquotedblright{} (EtymDB, under English \emph{symptom}), from \emph{sympíptein} \textquotedblleft{}to befall\textquotedblright{}.
\end{cousinweb}

\begin{grammarlens}[title={What you've built}]
One more word for telling, explaining and describing at B1.
\end{grammarlens}

\subsection*{Your turn}
\begin{itemize}
\raggedright
  \item \emph{Say it:} \emph{il sintomo}
  \item \emph{Say it:} \emph{il sintomo}, once more
  \item \emph{Say it:} a sentence of your own with \emph{il sintomo}
  \item \emph{Recall:} say the Italian for anaesthesia, then the Italian for to fall ill, then say \emph{il sintomo} again
\end{itemize}

\begin{tcolorbox}[breakable,colback=teal!4,colframe=teal!35!black,title={Before you move on}]
What does il sintomo mean? (\textquotedblleft{}Symptom\textquotedblright{}.) Say it once more.
\end{tcolorbox}
\section[la diagnosi]{la diagnosi — diagnosis}

\label{lesson:IT-C489-diagnosi}



\begin{quote}
Before the new one: say the Italian for to fall ill, then the Italian for symptom.
\end{quote}

\subsection*{You'll want to know: la diagnosi}
\textbf{la diagnosi} — \textquotedblleft{}diagnosis\textquotedblright{}.

\begin{quote}
La diagnosi è chiara: è solo un raffreddore.
\end{quote}

\begin{quote}
\emph{The diagnosis is clear: it is only a cold.}
\end{quote}

The same form in the plural: \emph{le diagnosi}.

\begin{cousinweb}
\emph{diagnosi} is a learned borrowing, through Latin \emph{diagnōsis}, from Greek \emph{diágnōsis} (EtymDB, under English \emph{diagnosis}), from \emph{diagign\'{\={o}}skein} \textquotedblleft{}to discern, to tell apart\textquotedblright{}.
\end{cousinweb}

\begin{grammarlens}[title={What you've built}]
One more word for telling, explaining and describing at B1.
\end{grammarlens}

\subsection*{Your turn}
\begin{itemize}
\raggedright
  \item \emph{Say it:} \emph{la diagnosi}
  \item \emph{Say it:} \emph{la diagnosi}, once more
  \item \emph{Say it:} a sentence of your own with \emph{la diagnosi}
  \item \emph{Recall:} say the Italian for to fall ill, then the Italian for symptom, then say \emph{la diagnosi} again
\end{itemize}

\begin{tcolorbox}[breakable,colback=teal!4,colframe=teal!35!black,title={Before you move on}]
What does la diagnosi mean? (\textquotedblleft{}Diagnosis\textquotedblright{}.) Say it once more.
\end{tcolorbox}
\section[l'infezione]{l'infezione — infection}

\label{lesson:IT-C489-infezione}



\begin{quote}
Before the new one: say the Italian for symptom, then the Italian for diagnosis.
\end{quote}

\subsection*{You'll want to know: l'infezione}
\textbf{l'infezione} — \textquotedblleft{}infection\textquotedblright{}.

\begin{quote}
Il medico mi ha dato un antibiotico per l'infezione.
\end{quote}

\begin{quote}
\emph{The doctor gave me an antibiotic for the infection.}
\end{quote}

\emph{Fare infezione} — to get infected.

\begin{cousinweb}
\emph{infezione} is a learned borrowing from Late Latin \emph{infectiō} \textquotedblleft{}a dyeing, a staining\textquotedblright{} (EtymDB), from \emph{inficere} \textquotedblleft{}to dip, to taint\textquotedblright{}, built on \emph{facere} \textquotedblleft{}to make\textquotedblright{}.
\end{cousinweb}

\begin{grammarlens}[title={What you've built}]
One more word for telling, explaining and describing at B1.
\end{grammarlens}

\subsection*{Your turn}
\begin{itemize}
\raggedright
  \item \emph{Say it:} \emph{l'infezione}
  \item \emph{Say it:} \emph{l'infezione}, once more
  \item \emph{Say it:} a sentence of your own with \emph{l'infezione}
  \item \emph{Recall:} say the Italian for symptom, then the Italian for diagnosis, then say \emph{l'infezione} again
\end{itemize}

\begin{tcolorbox}[breakable,colback=teal!4,colframe=teal!35!black,title={Before you move on}]
What does l'infezione mean? (\textquotedblleft{}Infection\textquotedblright{}.) Say it once more.
\end{tcolorbox}
\section[guarire]{guarire — to recover, to get better; to heal}

\label{lesson:IT-C489-guarire}



\begin{quote}
Before the new one: say the Italian for diagnosis, then the Italian for infection.
\end{quote}

\subsection*{You'll want to know: guarire}
\textbf{guarire} — \textquotedblleft{}to recover, to get better; to heal\textquotedblright{}.

\begin{quote}
Sono guarito in pochi giorni.
\end{quote}

\begin{quote}
\emph{I got better in a few days.}
\end{quote}

With \emph{essere} when you get better (\emph{sono guarito}), with \emph{avere} when you heal someone (\emph{il medico mi ha guarito}).

\begin{cousinweb}
\emph{guarire} comes from a Germanic verb, Frankish \emph{warjan} \textquotedblleft{}to protect, to defend\textquotedblright{} (EtymDB): first to guard, then to cure.
\end{cousinweb}

\begin{grammarlens}[title={What you've built}]
One more word for telling, explaining and describing at B1.
\end{grammarlens}

\subsection*{Your turn}
\begin{itemize}
\raggedright
  \item \emph{Say it:} \emph{guarire}
  \item \emph{Say it:} \emph{guarire}, once more
  \item \emph{Say it:} a sentence of your own with \emph{guarire}
  \item \emph{Recall:} say the Italian for diagnosis, then the Italian for infection, then say \emph{guarire} again
\end{itemize}

\begin{tcolorbox}[breakable,colback=teal!4,colframe=teal!35!black,title={Before you move on}]
What does guarire mean? (\textquotedblleft{}To recover, to get better; to heal\textquotedblright{}.) Say it once more.
\end{tcolorbox}
